\section{Тест производительности}

\subsection{Методика}

Тест сравнивает построение обобщённого суффиксного дерева и поиск ответа с классическим
динамическим программированием для задачи наибольшей общей подстроки. DP хранит таблицу
размера $(|S_1|+1)(|S_2|+1)$ и работает за $O(|S_1|\cdot|S_2|)$ по времени и памяти.
Суффиксное дерево строится за $O(|S_1|+|S_2|)$, после чего выполняется один DFS.

На малых размерах бенчмарк сравнивает полные множества ответов двух алгоритмов. На больших
размерах DP пропускается, поскольку таблица становится слишком большой.

\subsection{Результаты}

\begin{table}[!ht]
\centering
\begin{tabular}{|r|r|r|r|}
\hline
\rowcolor{lightgray}
$|S_1|=|S_2|$ & Суффиксное дерево, мкс & DP, мкс & Проверка \\
\hline
1\,000 & 996 & 4\,847 & yes \\
\hline
2\,000 & 1\,329 & 24\,298 & yes \\
\hline
4\,000 & 3\,482 & 106\,665 & yes \\
\hline
\end{tabular}
\caption{Сравнение с динамическим программированием, алфавит из 4 букв}
\end{table}

\begin{table}[!ht]
\centering
\begin{tabular}{|r|r|r|}
\hline
\rowcolor{lightgray}
$|S_1|=|S_2|$ & Суффиксное дерево, мкс & Длина ответа \\
\hline
10\,000 & 12\,783 & 12 \\
\hline
50\,000 & 208\,406 & 15 \\
\hline
100\,000 & 339\,247 & 17 \\
\hline
\end{tabular}
\caption{Большие тесты, DP пропущен}
\end{table}

\begin{alltt}
# python3 lab5/benchmark/generator.py 1000 1000 4 42
# ./build/lab5/lab5_benchmark < tests.txt
Suffix tree LCS time: 996us
LCS length: 10
LCS count: 1
DP LCS time: 4847us
DP LCS length: 10
DP LCS count: 1
Results equal: yes

# python3 lab5/benchmark/generator.py 4000 4000 4 42
# ./build/lab5/lab5_benchmark < tests.txt
Suffix tree LCS time: 3482us
LCS length: 12
LCS count: 1
DP LCS time: 106665us
DP LCS length: 12
DP LCS count: 1
Results equal: yes

# python3 lab5/benchmark/generator.py 100000 100000 4 42
# ./build/lab5/lab5_benchmark < tests.txt
Suffix tree LCS time: 339247us
LCS length: 17
LCS count: 1
DP LCS skipped: input is too large
\end{alltt}

\subsection{Анализ результатов}

На малых тестах суффиксное дерево уже быстрее DP: при длине строк 4\,000 выигрыш
составляет примерно \textbf{30 раз}. При дальнейшем росте размера DP становится
непрактичным из-за квадратичной таблицы, тогда как суффиксное дерево продолжает
работать в пределах ограничений. Бенчмарк также проверяет корректность: для малых
тестов множества найденных подстрок у дерева и DP совпадают.

\pagebreak
